\section{ĐỀ XUẤT QUY TRÌNH TIỀN XỬ LÝ}
\subsection{Tiêu chí lựa chọn}
Nhóm lựa chọn quy trình theo mục tiêu nhận dạng đủ 43 lớp biển báo, nên dùng Macro-F1 trên validation làm tiêu chí chính. Accuracy và Weighted-F1 được dùng để kiểm tra thêm. PSNR và SSIM giúp giải thích ảnh đã thay đổi như thế nào, nhưng không thay thế kết quả nhận dạng. Trong nhóm tăng sáng, diffusion được chọn trước. Sau khi khảo sát thêm khử nhiễu và thứ tự kết hợp, Gaussian rồi Gamma đạt Macro-F1 validation cao nhất là 77,77\% và được chọn để đánh giá test.

Quy trình được chọn không nhất thiết đạt điểm cao nhất ở mọi mức test. Chẳng hạn, CLAHE rồi Gaussian nhỉnh hơn ở mức nặng. Nếu nhóm đổi lựa chọn sau khi đọc kết quả này, tập test sẽ không còn giữ đúng vai trò kiểm tra cuối cùng. Nếu mục tiêu ứng dụng có thêm giới hạn về độ trễ, nhóm cũng cần đặt tiêu chí đó trước và lựa chọn bằng validation.

\subsection{Quy trình được lựa chọn}
\par\begin{figure}[!htbp]\centering\setlength{\fboxsep}{7pt}
\fbox{\parbox{0.88\linewidth}{\centering
\begin{minipage}[t]{0.47\linewidth}\centering
\textbf{A. Đánh giá mô phỏng}\par
Ảnh gốc GTSRB + ROI\\[3pt]$\downarrow$\\[3pt]
Crop + lề 5\%, RGB $64\times64$\\[3pt]$\downarrow$\\[3pt]
Gamma làm tối, Poisson,\\Gaussian và clip
\end{minipage}\hfill
\begin{minipage}[t]{0.47\linewidth}\centering
\textbf{B. Luồng áp dụng}\par
ROI thiếu sáng đã có sẵn\\[3pt]$\downarrow$\\[3pt]
Crop + lề, RGB $64\times64$\\[3pt]$\downarrow$\\[3pt]
Giữ ảnh thiếu sáng đầu vào\\(không sinh suy giảm)
\end{minipage}\par\vspace{5pt}
$\searrow\qquad\qquad\swarrow$\\[3pt]
\textbf{Khối xử lý chung: RGB float32 trong $[0,1]$}\\[3pt]$\downarrow$\\[3pt]
Gaussian $5\times5$, $\sigma=0,8$, biên phản xạ\\[3pt]$\downarrow$\\[3pt]
Gamma tăng sáng $\alpha=0,8$ trên RGB float32\\[3pt]$\downarrow$\\[3pt]
Resize $224\times224$ + chuẩn hóa ImageNet một lần\\[3pt]$\downarrow$\\[3pt]
ResNet-18 cố định $\to$ một trong 43 lớp
}}\caption{Hai nguồn đầu vào dùng chung chuỗi được chọn trên validation. Chỉ nhánh A sinh suy giảm và dùng ảnh gốc để tính PSNR/SSIM ở $64\times64$; nhánh B mô tả cách áp dụng, chưa được đánh giá trên ảnh camera thực.}\label{fig:pipeline}\end{figure}\par

Khi áp dụng cho ảnh thiếu sáng đã có sẵn, nhóm chỉ chuẩn bị ảnh rồi đưa vào khối lọc và Gamma, không tạo thêm suy giảm. Luồng này giả định vùng biển báo đã được cung cấp và chưa có mô hình phát hiện vị trí. Nhãn và ảnh sáng tham chiếu không được đưa vào chuỗi xử lý. Chỉ bộ đánh giá dùng chúng để đo chất lượng và kiểm tra phân loại. Trong cài đặt, bộ lọc dùng \texttt{BORDER\_REFLECT\_101}, và chuỗi Gaussian--Gamma không chuyển ảnh sang 8 bit.

Nhóm kiểm tra ảnh đầu ra về kích thước, số kênh, miền giá trị và các giá trị không hữu hạn. Ba checkpoint nhận dạng được giữ nguyên, không cập nhật theo ảnh test. Dù điểm validation và F1 test được cải thiện, nhóm chưa xem đây là bảo đảm an toàn cho một hệ thống hỗ trợ lái xe thực tế.

\subsection{Đối chứng và tổ chức cài đặt}
Nhóm dùng Identity, Gamma, CLAHE, diffusion và Gaussian đơn lẻ để phân biệt lợi ích của các thành phần. Khi so hai thứ tự Gaussian--CLAHE, nhóm giữ clip limit 1, lưới $4\times4$ và cùng tham số Gaussian. Nếu đồng thời thay đổi cả độ mạnh và thứ tự, nhóm sẽ khó xác định chênh lệch đến từ yếu tố nào.

Mã tạo suy giảm nằm trong \path{src/degradation.py}, còn các bộ tăng cường nằm trong \path{src/enhancement.py}. Nhóm cài đặt chuỗi kết hợp tại \path{src/tools/processing_pipeline.py} và dùng \path{src/tools/processing_study.py} để chọn tham số, cố định giao thức đánh giá. Tệp \path{src/metrics.py} tính các chỉ số, còn \path{configs/} lưu cấu hình. Kết quả đi kèm tham số, cách chia tập, mã ảnh và hash checkpoint để người đọc có thể kiểm tra từng phép so sánh.

\subsection{Hướng đánh giá tiếp theo}
Nếu tiếp tục đề tài, nhóm sẽ ưu tiên bổ sung ảnh thiếu sáng thực có nhãn biển báo phù hợp. Nhóm cần xem các lỗi về nét số, đường viền và màu, sau đó mở rộng sang ảnh toàn cảnh có mô hình phát hiện biển báo. Khi không có ảnh sáng ghép cặp chính xác, nhóm không thể tính PSNR và SSIM theo cách của thí nghiệm tổng hợp. Nhóm cũng cần thử thêm kiến trúc nhận dạng, điều kiện chụp và lần huấn luyện độc lập. Nếu xác minh được track test, nhóm nên lấy mẫu lại theo track thay vì từng ảnh. Khoảng bootstrap hiện tại chưa tính được nguồn phụ thuộc này.
